\section{Discussion and open questions}\label{sec:discussion}

The empirical results of Section~\ref{sec:results} show that the
SMC${}^2$ filter and controller, composed in receding-horizon form,
produce a closed-loop policy that improves mean autonomic amplitude by
$+17\%$ over a sensible constant baseline on the FSA-v2 model, with no
fatigue-ceiling violations. This is encouraging, but the present
treatment is incomplete in three specific and important ways. Each of
the three subsections below identifies one outstanding analysis that
should be carried out before the model is used in a setting where its
properties matter --- for instance, in a real-time clinical setting
where mis-identification of a parameter could mean prescribing the wrong
training load.

\subsection{Outstanding analysis: Fisher information rank
identifiability}\label{sec:discussion-fim}

Remark~\ref{rem:fim} of Section~\ref{sec:fsa-remarks} introduced the
Fisher information matrix \eqref{eq:fim} and pointed out that the rank
check on $\Ical(\theta^*)$ has not been carried out for FSA-v2. We now
sketch what such an analysis would look like and what its consequences
might be.

\paragraph{Procedure.} Pick a representative truth point $\theta^*$ (the
parameter vector of Table~\ref{tab:fsa-truth} extended with the
observation-model coefficients). Generate $L$ independent simulated
trajectories of the model at $\theta^*$, each producing a record
$y^{(\ell)}_{1:T}$ of four-channel observations on the project's
fifteen-minute time grid. For each trajectory, evaluate the score
$s^{(\ell)}(\theta) = \nabla_\theta \log p(y^{(\ell)}_{1:T} \mid \theta)$
at $\theta = \theta^*$ by automatic differentiation through the
generative code (the JAX dynamics solver of
\texttt{version\_2/models/fsa\_high\_res/\_dynamics.py} together with the
observation model in \texttt{simulation.py}). Approximate the Fisher
information by the empirical second moment
\(
  \widehat\Ical(\theta^*) = L^{-1} \sum_\ell s^{(\ell)}(\theta^*)\,s^{(\ell)}(\theta^*)^{\!\top}
\)
and read off its singular-value decomposition. Repeat for several truth
points to check that the conclusions are not parameter-specific.

\paragraph{Possible findings.} If $\rank \widehat\Ical(\theta^*) = d$ and
the smallest singular value is bounded away from zero, the model is
locally identifiable at $\theta^*$ and inference is well-posed. If
instead the rank is deficient by $r$, the right-singular vectors
corresponding to the smallest $r$ singular values give the directions in
parameter space along which the likelihood is flat. Each such direction
should then be examined physically: it might indicate (i) a
collinearity between two physiological parameters that no four-channel
observation set could resolve, in which case one of them must be fixed
or tightly priored; (ii) a redundancy in the observation model itself
(for instance, that HR and stress carry overlapping information about
$B$ and $F$); or (iii) a genuine modelling ambiguity that should be
removed by reparameterisation. Consequences for the SMC${}^2$ filter
include the possibility of multimodal posteriors, slower mixing of the
HMC moves of Section~\ref{sec:smc2}, and degraded robustness of the
closed-loop policy in the directions that the filter cannot constrain.

\subsection{Outstanding analysis: Lyapunov stability}\label{sec:discussion-lyap}

Remark~\ref{rem:lyap} flagged the Lyapunov analysis as an outstanding
piece of work in two parts: stability of the deterministic skeleton, and
stochastic stability of the full SDE.

\paragraph{Deterministic skeleton.} For a constant control
$\Phi = \Phi_{\mathrm{const}}$, the steady-state $A^*=0$ branch satisfies
\(
  B^*  = \kappa_B(1+\varepsilon_A \cdot 0)\Phi_{\mathrm{const}}\,\tau_B = \kappa_B \tau_B \Phi_{\mathrm{const}}
\)
and
\(
  F^* = \kappa_F\,\tau_F\,\Phi_{\mathrm{const}}.
\)
At this branch the Stuart--Landau growth rate is
$\mu(B^*,F^*) = \mu_0 + \mu_B B^* - \mu_F F^* - \mu_{FF} (F^*)^2$, and
the deterministic $A$-equation is locally $\dot A = \mu(B^*,F^*)\,A$.
The $A=0$ branch is therefore stable when $\mu(B^*,F^*) < 0$ and unstable
when $\mu(B^*,F^*) > 0$. In the latter regime an additional stable
limit cycle of radius $\sqrt{\mu(B^*,F^*)/\eta}$ appears. The complete
analysis would chart the bifurcation locus
$\{\Phi_{\mathrm{const}} : \mu(B^*(\Phi_{\mathrm{const}}),F^*(\Phi_{\mathrm{const}})) = 0\}$
and identify the critical training stimulus at which the autonomic
amplitude switches from quiescent to oscillatory.

\paragraph{Stochastic stability.} The natural candidate Lyapunov function
on $\R^3_+$ is the quadratic
\(
  V(B,F,A) = a B^2 + b F^2 + c A^2,
\)
with $a,b,c > 0$ to be chosen. Under It\^o calculus,
\begin{align*}
  \Lcal V(B,F,A)
    &= 2a B \bigl[\kappa_B(1+\varepsilon_A A)\Phi - B/\tau_B\bigr]
       + a\,\sigma_B^2\,B(1-B) \\
    &\quad
       + 2b F \bigl[\kappa_F \Phi - (1+\lambda_A A)/\tau_F \cdot F\bigr]
       + b\,\sigma_F^2\,F \\
    &\quad
       + 2c A \bigl[\mu(B,F)\,A - \eta A^3\bigr]
       + c\,\sigma_A^2\,A.
\end{align*}
The terms $-2a B^2/\tau_B$, $-2b F^2 (1+\lambda_A A)/\tau_F$ and
$-2c \eta A^4$ are dissipative; the remaining terms grow at most
linearly or quadratically in the state. For sufficiently large $a,b,c$
and appropriate operating range one would hope to establish a
Foster--Lyapunov drift inequality $\Lcal V \le -\alpha V + c'$ on
$\{V \le R\}$ for some $R$, which would imply exponential ergodicity and
the existence of a unique invariant distribution. The argument is
delicate near the Stuart--Landau bifurcation, where the $A$-direction is
weakly stabilising or weakly unstabilising; a piecewise Lyapunov function
that switches between regimes may be required. None of this has yet been
carried out.

\subsection{Outstanding analysis: LQG / Riccati exact-approximate
alternative}\label{sec:discussion-lqg}

Section~\ref{sec:mpc-lqg} sketched a linear--quadratic--Gaussian
controller for FSA-v2 under Assumptions~\ref{ass:lqg-A1}--\ref{ass:lqg-A3}.
None of the corresponding code is currently in the repository. We
identify three concrete pieces of work that together would deliver the
LQG companion to the SMC${}^2$ controller:
\begin{enumerate}
  \item Implement automatic linearisation of the drift around a
        user-supplied operating point, using the existing JAX dynamics in
        \texttt{version\_2/models/fsa\_high\_res/\_dynamics.py}; expose
        the time-varying $A_{\mathrm{lin}}(t),B_{\mathrm{lin}}(t)$ of
        \eqref{eq:lqg-AB} on the same fifteen-minute time grid used by
        the rest of the project.
  \item Implement a backward-integration solver for the differential
        Riccati equation \eqref{eq:riccati} on the same time grid.
  \item Build an \texttt{LQGController} class that emits the feedback
        gain $K(t)$ of \eqref{eq:lqg-feedback} and a corresponding
        $\Phi$ schedule on the time grid.
\end{enumerate}
Two scientific questions then become accessible. \emph{Does the LQG
schedule, evaluated as an open-loop control on the FSA-v2 model, achieve
mean amplitude comparable to the SMC${}^2$ schedule on the $84$-day
benchmark of Section~\ref{sec:results}?} If yes, the non-linear
SMC${}^2$ machinery may be unnecessary in steady-state regimes; if no,
the gap quantifies the value of capturing the cubic Stuart--Landau
non-linearity. \emph{Does using the LQG schedule as a warm-start for the
SMC${}^2$ control of Section~\ref{sec:control} reduce wall-clock time to
solution?} A reasonable hypothesis is that LQG provides an excellent
initial $\thetactrl$ in the linear regime and only marginal benefit
where the non-linearity matters; experimentally testing this is the
natural next step.

\subsection{Other open questions}\label{sec:discussion-other}

A few smaller points deserve mention but are not, in the authors' view,
on the same level of importance as the three analyses above.

\begin{itemize}
  \item \emph{Horizon length.} The Stage~E5 result uses a fourteen-day
        macrocycle. The Stage~D open-loop result uses eighty-four days.
        It is not yet known how the closed-loop gain scales with horizon
        length; a horizon sweep is straightforward to set up.
  \item \emph{Posterior shape vs.\ posterior mean.} The controller of
        Section~\ref{sec:mpc-pipeline} uses only the posterior mean
        $\bar\theta_{\mathrm{ctrl}}$. The flat-cost-surface argument of
        Section~\ref{sec:mpc-robustness} explains why this is sometimes
        sufficient, but in regimes where the cost surface is curved the
        full posterior carries useful information. A risk-sensitive
        variant that minimises a quantile of the cost rather than the
        mean would be the natural extension.
  \item \emph{Parameter drift across windows.} The Bures--Wasserstein
        bridge of Section~\ref{sec:smc2} produces a sensible
        Gaussian-fit warm-start, but does not protect against
        \emph{slow} parameter drift in the truth (e.g.\ the athlete's
        time constants change with training age). A more principled
        approach would augment the static parameters with a slow
        random-walk component and let the filter track the drift
        explicitly.
\end{itemize}

\medskip
\noindent
The empirical successes reported in Section~\ref{sec:results} should be
read alongside the three outstanding analyses above. The model is
empirically useful; it is not yet mathematically complete.
